\section{Introduction}
\label{sec:intro}

Finding a math problem that is the same problem written differently matters for spotting duplicates and for handing solved problems to a reasoning model. Many benchmarks now test retrieval over math and other reasoning problems~\citep{su2025bright,alshammari2026mathnet,georgiev2026saber,ye2026meld}, and more of them are built by an LLM: it writes the queries or the gold documents and decides what counts as relevant. When every gold document comes out of one fixed LLM procedure, a model can score well by learning what that procedure produces instead of the skill the benchmark says it tests. We call such gains \emph{recipe-matching}. Unlike an ordinary shortcut, a quirk of one dataset, this one is built into how the benchmark is made, so it carries over when a different LLM runs the same procedure and needs no access to the exact prompt (Section~\ref{sec:related}).

MathNet-Retrieve~\citep{alshammari2026mathnet} is a good place to study this. Its three tiers use the same corpus and the same queries, and differ only in how heavily the gold document is disguised. Every gold is a Gemini-3-flash paraphrase of its query, together with three near-misses, small edits of the query that change the problem, and it is kept only when Gemini-3-flash and GPT-5 both judge it equivalent. The hard tier is presented as deep mathematical equivalence, yet at rank~$1$ nobody solves it: every model in the authors' evaluation scores almost zero R@1, and no public model we tested exceeds $0.01$. That floor is built into the tier's design rather than a sign of a hard skill: the near-misses are tiny edits of the query, so any embedder that goes by surface wording ranks them above the heavily disguised gold, and for existing models R@5 and R@10 say more (Appendix~\ref{app:results}).

To find out what that tier rewards, we ran a controlled experiment on Qwen3-Embedding-0.6B~\citep{qwen3embedding}: two training sets of exactly $6{,}145$ rows built from the same source problems and trained under the same settings, so the only difference is the training file. In the \emph{verified arm}, a computer algebra system proves each positive equivalent to its source and refutes each minimal-edit negative with a counterexample. In the \emph{recipe arm}, the pairs are generated with the benchmark's own Appendix-F prompt and kept or dropped by an LLM judge alone, which is how the benchmark itself was built, except that our generator is Qwen3-32B rather than Gemini-3-flash and one judge decides where the benchmark requires two to agree. Both arms carry the same item-level contamination~\citep{dekoninck2024constat,yao2024crosslingual}, about a third of their rows anchoring on a benchmark query; it accounts for about one point of the gap, which survives two corrected-gate retrainings (Section~\ref{sec:setup}, Appendix~\ref{app:contamination}).

The recipe arm beats the verified arm on every tier (Table~\ref{tab:controlled}): over eight seeds it leads by $45.33$ easy-tier R@1 points and reaches $9.42$ on the hard tier, where the verified arm sits at zero. A third model, the \emph{recipe-free reference}, trains on LLM rewrites from a prompt unrelated to the benchmark plus the verified arm's negatives; it recovers $29.85$ of the $45.33$ points and matches the recipe arm on the hard tier, a second unrelated prompt recovers $22.15$, and back-translated paraphrases, which no LLM wrote, recover only $1.76$ with the same negatives. Half to two thirds of the gap is therefore bought by LLM-written pairs of any kind, and $15.47$ to $24.59$ points, depending on the comparison, need the benchmark's own prompt (Section~\ref{sec:gaming}). A ladder of rewriter prompts and a factorial over positives and negatives (Section~\ref{sec:mechanism}) show that the easy tier rewards LLM-written text in general, the \emph{genre} of LLM rewrites, whichever prompt wrote them. The hard tier rewards the \emph{pair structure} the recipe produces, a deep rewrite set against a minimal-edit near-miss: rewrites alone or near-misses alone score nothing, and restatements from an unrelated prompt set against the recipe's own near-misses reach $26.97$ over three seeds, nearly three times the recipe arm.

To see whether the lead means anything real, we test on duplicates no LLM wrote: the same problem printed in two languages, and the same problem reprinted in one. On the cross-language duplicates the recipe arm's easy-tier lead over the verified arm shrinks from $45.33$ to under ten points; the drop from benchmark to real data is $35.59$ points, and $27.24$ of it is recipe-specific, since the reference, with no recipe wording, loses far less. On same-language reprints the recipe arm holds no measurable lead over the verified arm and is within noise of the reference (Section~\ref{sec:ood}).

We tried the same attack on two other benchmarks. MELD~\citep{ye2026meld} never published the prompt behind its evaluation set; we reconstructed the procedure from its description and ran it through a different vendor's LLM, and training a $0.6$B model on the result more than triples its R@1 on MELD's pairs. That model also improves, rather than worsens, on cross-language duplicates, so most of its gain is the task itself being learned. SABER-Math~\citep{georgiev2026saber} has organic documents, the attack we registered against them fails, and the only movement is a small gain, robust to a matched budget, from training on the LLM-written summaries its pipeline mines pairs from (Section~\ref{sec:ood}).

The model trained on rewrites from an unrelated prompt with no negatives is at the top on cross-language duplicates, above the untrained base, yet scores zero on the hard tier; one edit to its training, attaching verified negatives, lifts it off that floor and costs it $15.79$ cross-language points. We release the evaluations, near-miss tests, three trained models and recommendations to benchmark builders.
